\subsection{Two edges and comparison with the QWZ measurement}

\begin{corollary}[Two-edge quarter twist]\label{cor:twoedge}
For two independent chiral edge branches, one at each boundary of a
cylinder, the direct-sum polarizations satisfy
\[
 \norm{P_0^{\mathrm{edge}}-P_\alpha^{\mathrm{edge}}}_2^2
 =\frac{4\sin^2(\pi\alpha)}{\pi^2}\log N+O_\alpha(1).
\tag{5.9}\label{eq:two-edge-log}
\]
For the adjacent $\mu_4$ sectors, $\alpha=1/4$, hence
\[
 \norm{P_0^{\mathrm{edge}}-P_{1/4}^{\mathrm{edge}}}_2^2
 =\frac2{\pi^2}\log N+O(1),
\qquad
\frac2{\pi^2}=0.2026423672846755\ldots .
\tag{5.10}\label{eq:quarter-coefficient}
\]
The difference from the measured slope $0.203$ is
$0.0003576327\ldots$, or approximately $0.1762\%$ of $0.203$.
\end{corollary}

\begin{proof}
Hilbert--Schmidt norms square add under orthogonal direct sums.  The two
edges have opposite propagation orientation, but the overlap formula
depends on the squared modulus and is unchanged.  Doubling
\eqref{eq:one-edge-log} proves \eqref{eq:two-edge-log}.  Since
$\sin^2(\pi/4)=1/2$, the coefficient becomes $2/\pi^2$.  The numerical
comparison is direct arithmetic.
\end{proof}

\begin{lemma}[Stability under a uniformly Hilbert--Schmidt remainder]
\label{lem:stability}
Suppose finite-volume differences $D_N$ admit
\[
 D_N=D_N^{\mathrm{edge}}+R_N,\qquad
 \sup_N\norm{R_N}_2<\infty,
\tag{5.11}\label{eq:remainder}
\]
and
\[
 \norm{D_N^{\mathrm{edge}}}_2^2=c\log N+O(1)
\quad(c>0).
\]
Then
\[
 \norm{D_N}_2^2=c\log N+O(\sqrt{\log N}),
\tag{5.12}\label{eq:stable-log}
\]
so $\norm{D_N}_2\to\infty$ and the leading logarithmic coefficient is
still $c$.
\end{lemma}

\begin{proof}
Expand the square:
\[
 \norm{D_N}_2^2
 =\norm{D_N^{\mathrm{edge}}}_2^2
 +2\operatorname{Re}\Tr((D_N^{\mathrm{edge}})^*R_N)
 +\norm{R_N}_2^2.
\]
Hilbert--Schmidt Cauchy--Schwarz gives
\[
 \left|\Tr((D_N^{\mathrm{edge}})^*R_N)\right|
 \le\norm{D_N^{\mathrm{edge}}}_2\norm{R_N}_2
 =O(\sqrt{\log N}).
\]
The last term is $O(1)$ by \eqref{eq:remainder}.  This proves
\eqref{eq:stable-log}.
\end{proof}

\begin{theorem}[Global non-implementability in the proved edge class]
\label{thm:nonimpl}
Let $P_0$ and $P_{1/4}$ be the infinite-volume pure quasi-free
polarizations of either the exact two-edge model of
Corollary~\ref{cor:twoedge}, or a lattice collar for which the
uniform-remainder decomposition \eqref{eq:remainder} holds along an
exhaustion.  Then the two quasi-free representations are not unitarily
equivalent.  Equivalently, no global Fock-space Bogoliubov implementer
can carry the neutral polarization to the quarter-twisted polarization.
\end{theorem}

\begin{proof}
For the exact edge model, Theorem~\ref{thm:oneedge} and
Corollary~\ref{cor:twoedge} show that the nonnegative finite-rank
truncations of $\norm{P_0-P_{1/4}}_2^2$ tend to infinity.  Therefore
$P_0-P_{1/4}\notin\mathcal L^2$.  Under
\eqref{eq:remainder}, Lemma~\ref{lem:stability} gives the same conclusion.

Assume for contradiction that a global Bogoliubov implementer exists.
Let $R$ be the underlying one-particle orthogonal transformation, so
$P_{1/4}=RP_0R^*$.  The Shale--Stinespring criterion,
Theorem~\ref{thm:ss}, would imply
\[
 [P_0,R]\in\mathcal L^2.
\]
But
\[
 [P_0,R]R^*=P_0-RP_0R^*=P_0-P_{1/4},
\]
and multiplication by $R^*$ preserves the Hilbert--Schmidt class.
Hence $P_0-P_{1/4}$ would be Hilbert--Schmidt, contradicting the
divergence just proved.  The pure-state Araki criterion quoted after
Theorem~\ref{thm:ps} gives the equivalent representation-theoretic
formulation.
\end{proof}

\begin{remark}[Gapped control]
If instead the global projection difference converges in
Hilbert--Schmidt norm, Shale--Stinespring permits a global implementer.
The measured trivial-phase plateau near $0.546$ is consistent with this
alternative, but a finite plateau is not a proof of convergence.
\end{remark}

\section{Combined strongest theorem}

\begin{theorem}[Strongest theorem proved from the presently isolated
mathematical hypotheses]\label{thm:main}
Consider a sequence of finite quadratic fermion systems with four
$\mu_4$-twisted sectors and a nested family of fixed finite windows.
Assume the following precise hypotheses.
\begin{enumerate}[label=(H\arabic*)]
\item On each fixed window, the restricted covariances
$C_{N,w}^{(k)}$ converge in trace norm.  For the adjacent sectors
$k=0,1$, their limiting Araki--Wyss projections have distance strictly
less than one.  In particular, this holds when the two limiting window
covariances coincide.
\item The finite observable algebras form an equivariant unital tower,
and their clock unitaries obey
$u_N^4=\id$ and $\iota_N(u_N)=u_{N+1}$.
\item The grade-one class is labelled by the frozen lattice element
$[\lambda]$ with $[\lambda]^4=[0]$ and $[\lambda]^2=[v]$.
\item For the global topological conclusion, the infinite-volume
polarization difference is the two-edge quarter-twist model of
Section~\ref{sec:global}, up to a remainder uniformly bounded in
Hilbert--Schmidt norm along the exhaustion.
\end{enumerate}
Then:
\begin{enumerate}[label=\textup{(\alph*)}]
\item On every fixed window, the adjacent sector change has a canonical
Bogoliubov implementer in the purified standard representation, unique
up to phase, with the explicit bound
\[
\norm{[\widehat P_{C_{N,w}^{(0)}},R_{N,w}]}_2
\le2\sqrt{2\norm{C_{N,w}^{(0)}-C_{N,w}^{(1)}}_1}.
\]
\item Along every subsequence on which the phase is chosen coherently,
these implementers converge in operator norm, hence strongly on the
finite-particle core; estimate \eqref{eq:nparticle} holds for every
$n$-particle vector.
\item The inductive-limit extension is
\[
\A\rtimes\mathbb Z_4\cong\A\otimes C^*(\mathbb Z_4),
\]
the homogeneous subspaces have $\mathbb Z_4$ fusion,
$\Psi_\lambda^4\in\A$, and $\Psi_\lambda^2$ has the parity grade
$[v]$.  The coefficient expectation is the unique dual-invariant
expectation onto $\A$, and the group average is the unique
trace-preserving expectation onto $\A^{\mathbb Z_4}$.
\item In the topological two-edge class no global Bogoliubov implementer
exists.  The obstruction has the exact leading law
\[
\norm{P_0-P_{1/4}}_2^2
=\frac2{\pi^2}\log N+O(\sqrt{\log N}).
\]
\end{enumerate}
\end{theorem}

\begin{proof}
Parts (a) and (b) are Theorem~\ref{thm:local}, with all hypotheses
verified in its proof.  Part (c) is Theorem~\ref{thm:crossed}; hypothesis
(H3) supplies only the identification of the abstract grades with the
frozen lattice labels.  Part (d) follows from
Corollary~\ref{cor:twoedge}, Lemma~\ref{lem:stability}, and
Theorem~\ref{thm:nonimpl}.  No further implication is used.
\end{proof}

\section{What this theorem does and does not close}

\subsection{Relation to G1--G4}

\begin{description}[style=nextline]
\item[G1: convergence on the finite-particle core with energy bounds.]
Theorem~\ref{thm:local} proves strong convergence on each
\emph{fixed purified finite window}.  It does not prove convergence of
a nontrivial charged intertwiner on one common infinite-volume edge Fock
space.  The measured inequalities with $C=2.86$ concern only
$n=1,\ldots,4$ and finite $N_x$; they do not imply
\[
\norm{H_{\mathrm{edge}}^n\Psi_\lambda\Omega}
\le C_0^{n+1}(n!)^\beta\qquad\text{for every }n\ge0.
\]
Thus G1 remains open in its contract meaning.

\item[G2: locality relative to the bilinear net.]
No commutator estimate at spacelike separation, no Haag--Kastler
localisation, and no relative-locality theorem follows from covariance
convergence.  G2 is untouched.

\item[G3: braiding and monodromy.]
The exact arithmetic $h_\lambda=1$ is the correct trivial
self-monodromy precondition.  The measured Berry phase $0$ with spectral
flow $1$ is consistent with it.  Neither fact computes the DHR braiding
operator or proves the relative exchange relation of a charged field.
G3 is therefore not closed.

\item[G4: Q-system closure.]
Section~\ref{sec:tower} proves the $\mathbb Z_4$-graded crossed-product
algebra, its naturality in the finite tower, and the precisely normalised
expectations.  This is a complete $C^*$-algebraic statement for the inner
clock tower.  It does not identify that tower with the local
simple-current Q-system of the $\mathfrak{so}(16)_1$ conformal net.
Hence it is the exact finite/inductive algebraic shadow of G4, not the
net-level G4 demanded by the contract.
\end{description}

\subsection{The MMST/Osborne--Stottmeister leg}

T.~J.~Osborne and A.~Stottmeister,
\emph{Conformal field theory from lattice fermions},
Communications in Mathematical Physics \textbf{398} (2023), 219--283,
\href{https://arxiv.org/abs/2107.13834}{arXiv:2107.13834},
prove strong convergence of the lattice fermion-bilinear currents on a
finite-particle core (Theorem~5.1 for the displayed $u(1)$ current, with
the nonabelian bilinear extension by the Cauchy--Schwarz reduction around
their Eq.~(229)), convergence of the Koo--Saleur/Virasoro generators
(Theorems~4.11 and 4.16), and correlation-function control
(Theorems~6.1 and 6.3).  Their error estimates around Eqs.~(205) and
(207) are polynomial in lattice spacing.

Those theorems concern observables that are fermion bilinears.  The
$128$ weight-one spinor currents needed for
$\mathfrak{so}(16)_1\to(E_8)_1$ are not fermion bilinears, and the twist
field is outside that proved class.  Therefore the statement additionally
needed here is not another application of Shale--Stinespring.  It is an
identification theorem saying that one specifically normalised limit of
the lattice twist intertwiners is a relatively local DHR charged field
whose $\mathbb Z_4$ Q-system extension is the holomorphic $(E_8)_1$ net.
That assertion belongs to conformal-net/simple-current extension theory,
not to the quasi-free category: a spin field changes sectors and is not a
polynomial fermion bilinear.

\subsection{Status of the measured QWZ logarithm}

The coefficient $0.203$ has a sharp analytic explanation:
each chiral boundary contributes
\[
\frac{2\sin^2(\pi/4)}{\pi^2}=\frac1{\pi^2},
\]
and the two boundaries contribute $2/\pi^2=0.202642367\ldots$.
The measured intercept $3.605$ is nonuniversal and receives bulk,
ultraviolet, and finite-width contributions.  To turn the coefficient
match into a proof for the QWZ collar, one must establish the decomposition
\eqref{eq:remainder} uniformly in $N_x$.  That estimate is not contained
in the frozen probes.

\section{Remaining gap}

\begin{center}
\setlength{\fboxsep}{10pt}
\fcolorbox{red!65!black}{red!3}{
\begin{minipage}{0.91\textwidth}
\textbf{Remaining gap (exact statement).}
For the $M=1$, $N_y=8$ QWZ collar with quarter-twisted sectors, construct
isometric embeddings of all finite-volume one-particle spaces into one
edge Hilbert space and prove simultaneously:
\[
C_N^{(0)}-C_N^{(1)}
=D_{N,\mathrm{top}}^{(1/4)}
\oplus D_{N,\mathrm{bottom}}^{(1/4)}+R_N,
\qquad
\sup_N\norm{R_N}_{\mathrm{HS}}<\infty;
\]
and there exist phase-coherent lattice intertwiners
$\Psi_{\lambda,N}$ converging on one common finite-energy core to a
nonzero operator $\Psi_\lambda$ such that, for constants
$C_0<\infty$ and $\beta<\infty$ independent of $n$,
\[
\norm{H_{\mathrm{edge}}^n\Psi_\lambda\Omega}
\le C_0^{\,n+1}(n!)^\beta\quad(n\ge0),
\]
$\Psi_\lambda$ is relatively local with respect to the
$\mathfrak{so}(16)_1$ bilinear net, its DHR exchange operators realise
the order-four class with $h_\lambda=1$ and
$[\lambda]^2=[v]$, and the local extension generated by
$\{\mathcal A_{\mathfrak{so}(16)_1},\Psi_\lambda\}$ is isomorphic, as a
conformal net with vacuum and stress tensor, to the holomorphic
$(E_8)_1$ net.

\medskip
The first line upgrades the finite logarithmic fit to the proved Shale
obstruction for the actual collar.  The remaining lines are the
MMST/Osborne--Stottmeister extension beyond fermion bilinears and supply
G1--G4 at net level.  Neither assertion follows from the measured
windowed Hilbert--Schmidt Cauchy rate, three phase-fixed matrix elements,
or four finite energy moments.
\end{minipage}}
\end{center}

\section*{Conclusion}
\addcontentsline{toc}{section}{Conclusion}

The quasi-free conclusions that can be proved today are complete:
finite-window standard-form implementers with an explicit Shale bound,
their coherent strong convergence, the natural inner-$\mathbb Z_4$
crossed-product tower with correctly qualified unique expectations, and
the universal two-edge quarter-twist obstruction with coefficient
$2/\pi^2$.  The data strongly match this structure, but they do not yet
construct the net-level field $\Psi_\lambda$.  The remaining box states
exactly the theorem that would be required to do so.

\begin{thebibliography}{9}

\bibitem{Araki1970}
H.~Araki,
\emph{On quasifree states of CAR and Bogoliubov automorphisms},
Publ. Res. Inst. Math. Sci. \textbf{6} (1970/71), 385--442.

\bibitem{OsborneStottmeister2023}
T.~J.~Osborne and A.~Stottmeister,
\emph{Conformal field theory from lattice fermions},
Commun. Math. Phys. \textbf{398} (2023), 219--283,
\href{https://arxiv.org/abs/2107.13834}{arXiv:2107.13834}.

\bibitem{PowersStormer1970}
R.~T.~Powers and E.~St{\o}rmer,
\emph{Free states of the canonical anticommutation relations},
Commun. Math. Phys. \textbf{16} (1970), 1--33.

\bibitem{Ruijsenaars1978}
S.~N.~M.~Ruijsenaars,
\emph{On Bogoliubov transformations. II. The general case},
Ann. Physics \textbf{116} (1978), 105--134.

\bibitem{ShaleStinespring1964}
D.~Shale and W.~F.~Stinespring,
\emph{States of the Clifford algebra},
Ann. of Math. (2) \textbf{80} (1964), 365--381.

\end{thebibliography}

\end{document}
